%\documentclass[english,9pt]{DLRK_v1.11}
\documentclass[german,10pt]{DLRK_v1.11}





% % Please provide the title of the paper hereafter (non-capitalized as it will be 
% % automatically capitalized in the paper, but let unchanged for the field "Title" 
% % of the PDF file metadata)
\newcommand{\PaperTitle}{Richtlinien zur Erstellung von Manuskripten zur Ver\"{o}ffentlichung in Kongressdokumentationen}



% % Please provide the names of the authors hereafter twice:
% %   A) first with the \ssymbol commands for symbols in the article 
% %      (use the same index for authors with the same affiliation, 
% %      otherwise start with 1 and increment them with each new author
% %      whose affiliation is not identical to one of the previous ones
% %   B) then repeat the list as shown but without the symbols (that 
% %      second list of authors will be used for the field "Author" of
% %      the PDF file metadata.
\newcommand{\PaperAuthors}{
  \texorpdfstring
    {A. Author1\ssymbol{1}, A. Author2\ssymbol{1}, A. B. Author3\ssymbol{2}} % first list
    {A. Author1, A. Author2, A. B. Author3}                                  % second list
}

% % Please provide the contact e-mail address (only one) to contact the 
% % authors of the paper
\newcommand{\ContaktEmailAddress}{corresponding@author.com}


% % Provide all the distinct affiliations of the authors hereafter
% % in the order of appearance in the list of authors
\newcommand{\PaperAuthorsAffiliations}{
\ssymbol{1}\ Name of the First Organization, Department Name, Address, City, Country\\
\ssymbol{2}\ Name of the Second Organization, Department Name, Address, City, Country\\
} % % CAUTION: Do not forget the end the last entry with two backslashes
  % % other wise the last affiliation might not be centered


\hypersetup{
    pdftitle={\PaperTitle},
    pdfauthor={\PaperAuthors},
}




\title{\PaperTitle}
\author{\PaperAuthors\vspace{0.5ex}\\\PaperAuthorsAffiliations}


%\pagestyle{empty}



\newcommand{\makeabstract}{\begin{onecolabstract}
Ein Abstract ist ein kurzer Abriss und eine knappe Inhaltsangabe. Gerade ungeduldige Leser lesen gerne den Abstract, um zu entscheiden, ob sich ein tiefer gehendes Lesen wirklich lohnt.
\end{onecolabstract}
}

\newcommand{\paperkeywords}{Keyword A; Keyword B; Keyword C}


\makenomenclature



% % Use these commands from the "nomencl" and "ifthen" packages to format
% % the nomenclature and subdivide it in adequate categories
% % NOTE THAT THE NAMES OF THE CATEGORIES IS NOT SET BY THE TEMPLATE BUT
% % BY THE AUTHORS
% % => NO AUTOMATIC TRANSLATION BETWEEN THE SUPPORTED LANGUAGES!!!!!!!!!!!
\renewcommand{\nomgroup}[1]{%
\ifthenelse{\equal{#1}{L}}{\pagebreak[2]\item[\textbf{Formelzeichen}]}{\linebreak
  \ifthenelse{\equal{#1}{S}}{\pagebreak[2]\item[\textbf{Indizes}]}{
    \ifthenelse{\equal{#1}{Z}}{\pagebreak[2]\item[\textbf{Abk\"urzungen}]}{
}}}}
% % Here is an example of entry for the nomenclature:
% %   \nomenclature[L_q]{$q$}{pitch rate\nomunit{deg/s}}
% %
% % where "L_q" is sorting key for the entry ("L" indicates that is a 
% % symbol and the following characters are used to automatically sort 
% % the symbols among each other)
% %
% % The second entry is the symbol itself (note that it is usually a 
% % good idea to define commands for the combined symbols (i.e. not just
% % one-letter symbols) and to use that command throughout the paper.
% % This way if you need to change the symbol (e.g. to avoid confusion
% % with some other symbol), you would only have to change its definition
% % in one location.
% % Note that for defining symbols that need to be used in the math 
% % environment (e.g. greek letter, symbols with sub- or superscripts)
% % the LaTeX command \ensuremath is very helpful.
% % 
% % Finally, the last argument in the description of the symbol. When
% % relevant the unit can be provided with the \nomunit command
% %
% % See the examples included in the included nomenclature file.
% % Note that for this nomencalture package to work properly, your
% % LaTeX compilation profile must include a call to "makeindex"
% % For windows users using TeXnicCenter this can be configured in 
% % the "Output Profile -> Postprocessing" pane with the following
% % configuration:
% %   - Executable: makeindex.exe
% %   - Arguments: "%tm".nlo -s nomencl.ist -o "%tm".nls



\begin{document}
\twocolumn[
\begin{@twocolumnfalse}
\maketitle
\vspace{0.3cm}

\makeabstract
\vspace{0.3cm}


\makepaperkeywords
\vspace{0.6cm}
\end{@twocolumnfalse}
]



% % In this example the items to be included in the nomenclature are defined
% % in a separate file "SymbolsAndAbbreviations.tex"
\input{SymbolsAndAbbreviations} 
% % Here to print the nomenclature: the argument is the width of the columns 
% % containing the symbols. This value can be adjusted (within some reasonnable
% % range) if needed.
\printnomenclature[1.7cm] 


\section{Gegenstand der Ver\"offentlichung}

Es können wissenschaftliche Arbeiten in digitaler Form veröffentlicht werden die auf Veranstaltungen der DGLR oder mit Beteiligung der DGLR vorgestellt wurden. Es werden nur Dokumente als Erstpublikation veröffentlicht.

\section{Form der Veröffentlichung}

Die Manuskripte werden in elektronischer Form (PDF/A-Datei) in öffentlichen Datennetzen (Internet) veröffentlicht. Jedes Dokument\begin{itemize}
\item enthält einen Hinweis auf die Veranstaltung und Metadaten mit Information über Autoren, Veranstaltung, Keywords und Rechte;
\item wird auf den Webseiten der DGLR der Öffentlichkeit kostenfrei zugänglich gemacht;
\item wird für die Langzeitarchivierung in ein PDF/A konvertiert und bei der Nationalbibliothek angemeldet;
\item erhält eine Uniform Resource Name (URN) zur eindeutigen Identifizierung und ermöglicht dauerhaft das Original aufzufinden.
\end{itemize}
Eine Veröffentlichung kann nur bei Einhaltung der Richtlinien garantiert werden.


\section{Formatierung}
\subsection{Seitenformat}
Papierform ist das Format DIN A4 (21\,cm $\times$ 29,7\,cm).
\subsubsection{Seitenränder}
Links und rechts jeweils 20\,mm, oben 25\,mm und unten 30\,mm.
\subsubsection{Spalten}
Vortragstitel, Autorenangaben, Zusammenfassung und Keywords sind auf der 1. Seite einspaltig anzugeben. Der Vortragstext, Abbildungen, Tabellen und Grafiken sind zweispaltig mit 5 mm zwischen den Spalten einzurichten.
\subsubsection{Kopf- und Fußzeilen}
Das Manuskript muss frei von Kopf und Fußzeilen sein und sollte keine Seitenzahlen enthalten. Für die Veröffentlichung werden automatisch Kopfzeilen mit
Informationen über Kongress und Dokument sowie
fortlaufende Seitennummerierung eingefügt.

\subsection{Textformatierung}
Der Text ist einzeilig, vorzugsweise in \enquote{Arial} mit einer Schriftgröße von mindestens 9 pt, zu erstellen.

\subsubsection{Titel der Arbeit und Autoreninformation}
Der Titel der Arbeit ist in Versalien (Großbuchstaben) einspaltig und zentriert zu schreiben. Darunter folgt mit Abstand und in einer separaten Zeile die Autoreninformation in der Form: Initial des Vornamens, Zuname, Name der Organisation mit Adresse und Staat. (bspw.: K.-H. Mustermann, Muster AG, 1234 Musterheim, Schweiz).

\subsubsection{Zusammenfassung}
Nach der Autoreninformation folgt als nächster Absatz einspaltig eine kurze Zusammenfassung der Arbeit. (Abstract). Darunter folgt der zweispaltige Text.



\subsubsection{Gliederungspunkte}

1. KAPITELÜBERSCHRIFTEN
(in Versalien und Fettdruck)
1.1. Abschnittsüberschriften
(in Groß- und Kleinschreibung und Fettdruck)
1.1.1. Unterabschnittsüberschriften
(in Groß- und Kleinschreibung und Fettdruck)

\subsubsection{Formelsymbole und Abkürzungen}
Werden entweder in einem eigenen Abschnitt (siehe zum Beispiel hier vor Abschnitt 1) oder im Text aufgeführt und erklärt.

\subsubsection{Abbildungen, Grafiken und Tabellen}
Sie können entweder im Text eingebunden oder als Anhang am Ende des Manuskriptes aufgeführt werden. Zwei Beispiele für Abbildungen und Tabellen sind hiernach in \figref{fig:longitudinal_controller} und \tabref{tab:abc} gegeben.

\begin{figure}
\centering
\includegraphics[width=\columnwidth]{./graphics/longitudinal_controller_vDLRKTemplate.pdf}
\caption{Flight path angle ($\gamma$) longitudinal control law}
\label{fig:longitudinal_controller}
\end{figure}


\begin{table}%
\centering
\begin{tabular}{lcr}\hline
a &b &c\\\hline
\end{tabular}
\caption{test table}
\label{tab:abc}
\end{table}


\subsubsection{Gleichungen}
Für die Nummerierung von Gleichungen werden Nummern in runden Klammern verwendet die linksbündig der Gleichung vorangestellt sind. Bitte vermeiden Sie die Schrift: Cambria Math.

\begin{equation}
F_{crit} = \gamma \frac{2\pi E t^2 \cos(\alpha)^2}{\sqrt{3(1-\mu^2)}}
\end{equation}


\subsubsection{Fu\ss noten}
Fußnoten sind im Text mit einem Asteriskus *) oder einem geeignetem Symbol kenntlich zu machen.

\subsubsection{Schrifttum}
Das Schrifttum soll am Schluss der Arbeit in der Reihenfolge der im Text nummerierten Zitate aufgelistet werden. Die fortlaufende Nummerierung ist in eckige Klammern zu setzen. Es wird Wert auf vollständige Angaben gelegt, d.h. bei Zeitschriftenaufsätzen sind Verfasser mit Anfangsbuchstaben des Vornamens, den vollständigen Nachnamen, der Titel der Arbeit, Zeitschrift mit Bandnummer, Jahreszahl sowie erste und letzte Seite der Arbeit, bei Büchern Verfasser und Herausgeber, Titel des Buches, Verlag, Erscheinungsort und Jahr, bei Proceedings-Bänden neben der Bezeichnung des Bandes, des Herausgebers und Erscheinungsjahr, der volle Titel der Arbeit sowie deren erste und letzte Seite zuzitieren.

Beispiele für unterschiedliche Typen von Dokumenten (Zeitschriftartikel, Konferenzbeitrag, Dissertation, Buch, usw.) sind hiernach zu finden. Beim \LaTeX{} template wird ein DLRK ``bibliographystyle'' verwendet, welcher das ``doi'' Feld erkennt und automatisch einen Hyperlink produziert. Dieses erfordert, dass nur die reine DOI Nummer im BibTeX Eintrag eingegeben wird: Zum Beispiel nur ``10.1007/s13272-017-0240-9'' für Ref.~\cite{FezSchFis2017}. Für Bücher erkennt das BibTeX Style erkennt auch das ``isbn'' Feld, z.\,B. siehe Ref.~\cite{zhou1996robust}.\nocite{PhD_Thesis_Koethe,SchFezNie2019,NASA_TM4511,ISO1151-1:1988}

\subsubsection{Anhang}
Bei Bedarf kann am Ende der Arbeit für Abbildungen, Grafiken und Tabellen ein Anhang angefügt werden.

\section{Datenformat}
Für die Publizierung der einzelnen Beiträge wird das Portable Document Format (PDF) verwendet. Wir bitten deshalb um Einreichung Ihres Beitrages als PDF-Datei, möglichst als PDF/A. Andere Dateiformate sind nur nach vorheriger Absprache zur Einreichung verwendbar. Das Dokument sollte mindestens 4 Seiten umfassen, 10 Seiten sind erwünscht, und nicht größer als 6 MB sein. Achtung: Von Ihnen eingereichte bereits fertige PDF-Dateien dürfen keine Sicherheitseinstellungen aufweisen!

\section{Ver\"offentlichungsrecht und Nut\-zungs\-rechte}

Die rechtliche Vereinbarung zwischen Autoren/Recht\-inhaber und der Deutschen Gesellschaft für Luft- und Raumfahrt, als veröffentlichender Verlag entnehmen Sie bitte der Rechteerklärung für die Veröffentlichung. Alle Dokumente werden als 'Copyright protected'
veröffentlicht. Alternativ haben die Autoren die Möglichkeit den Nutzern weitergehende
Nutzungsrechte (Creative Commons Lizenzen (CC)) einzuräumen. Mit der Rechteerklärung beauftragt der Unterzeichnende, stellvertretend für alle Autoren/Rechteinhaber, die DGLR, sein Dokument mit einer CC Lizenz oder als copyright protected zu veröffentlichen.

\section{Ver\"offentlichung in elektronischer Form}
Der Ablauf des Veröffentlichungsprozesses finden sie unter: \href{http://www.dglr.de/?id=3484}{\color{blue}{http://www.dglr.de/?id=3484}}. Die Autoren/Rechteinhaber behalten alle Rechte für eine weitere Verwertung des Dokuments. Für Fragen von Nutzern bzgl. Inhalt oder Rechte geben sie bitte eine Emailadresse an.


% If really necessary multiple contact e-mail addresses can be provided hereafter,
% but only one address is preferrable. It does not need to be the one of the first
% author.
\subsection*{\ContactSection}
\href{mailto:\ContaktEmailAddress}{\color{blue}\ContaktEmailAddress}


\bibliography{DLRK_sample}

% % Appendix sections, if needed, can be inserted herafter after 
% % uncommenting the following line
%\appendix



\end{document}